\section{Additional Implementation Details}
\label{app:implementation}

\subsection{Forecaster architecture and training}
\label{ssec:fm-spec}
\fmname{} is built on the GraphPFN graph-tabular foundation
backbone~\citep{HollmannEtAl2025TabPFN,
eremeev2025turningtabularfoundationmodels} with a LiMiX-16M tabular
transformer encoder ($12$ layers, $4$ attention heads, hidden dimension
$64$). The same backbone serves three heads (edge existence, binary;
edge weight, regression; and node-TVL change, regression), trained
jointly with a BCE$+$MSE composite loss. We fine-tune one checkpoint per
horizon $h\!\in\!\{1,4,8,12\}$ for $20$ epochs with AdamW (learning rate
$10^{-4}$ for heads, $10^{-5}$ for the backbone). Fine-tuning draws on
\corpusname{} snapshots from March~2020 to January~2025 ($104$ training
weeks, $12$ validation weeks for early stopping, $8$ held-out weeks for
internal evaluation); all data that influences weights or checkpoint
selection precedes the frozen 2025 leaderboard split
(\S\ref{app:bench-details}). Full-graph in-context inference fits a single
RTX~4090, and one weekly forecast across all four horizons completes in
well under a minute. These intrinsic forecaster numbers are kept out of
the main text because the paper's contribution is the operating point of
the full pipeline, not the FM's headline accuracy.

\subsection{Benchmark details}
\label{app:bench-details}

\paragraph{Six-axis schema.}
\begin{center}\small
\begin{tabularx}{\linewidth}{@{}llX@{}}
\toprule
ID                       & Capability tested                & Primary metrics \\
\midrule
\texttt{b1\_forecast}    & Temporal graph prediction        & PageRank/HHI/Gini MAE, trend consistency, Spearman; per $h\!\in\!\{1,4,8,12\}$ \\
\texttt{b2\_warning}     & Streaming anomaly detection      & Precision, recall, F1, lead time, alert stability; per event $\times$ budget \\
\texttt{b3\_calibration} & Predictive uncertainty           & PI coverage at $0.90$ target, PI width, ECE, CRPS \\
\texttt{b4\_stress}      & What-if scenario fidelity        & Loss MAE, distressed-count MAE, propagation depth MAE, target overlap@$k$ \\
\texttt{b5\_decision}    & Supervisory ticket quality       & Precision, recall@$k$, F1, target stability, judge score, FIR \\
\texttt{b6\_robustness}  & Data quality sensitivity         & Per-benchmark metrics under 5 degradation regimes; relative degradation \\
\bottomrule
\end{tabularx}
\end{center}
The six axes are non-overlapping (a system can excel at \texttt{b1} while
failing \texttt{b5}), decomposable along the four-layer framework so that
predictor (\texttt{b1}, \texttt{b3}, \texttt{b6}), monitor (\texttt{b2}),
scenario (\texttt{b4}), and decision (\texttt{b5}) ablations report
independently, and individually publishable so downstream work may extend a
single axis without re-running the suite.

\paragraph{Dataset split.}
\begin{center}\small
\begin{tabularx}{\linewidth}{@{}llrX@{}}
\toprule
Split      & Period             & Weeks       & Usage \\
\midrule
Train      & 2020-03--2024-06   & $\sim$222   & Forecaster training \\
Validation & 2024-07--2024-12   & $\sim$26    & Conformal calibration, threshold tuning \\
Test       & 2025-01--2025-08   & $\sim$33    & Frozen leaderboard \\
\bottomrule
\end{tabularx}
\end{center}

\paragraph{Historical warning windows.}
The \texttt{b2\_warning} event study is separate from the frozen 2025
leaderboard split. It evaluates the shared weighted-degree monitor around
Terra/Luna (2022-05-09), FTX (2022-11-07), and SVB/USDC (2023-03-10) with a
shorter $26$-week rolling baseline, chosen so each event window has enough
pre-event history. The 2025 leaderboard runs use the $42$-week monitor
baseline in Appendix~\ref{app:pipeline-math}.

\paragraph{Ground-truth definition.}
With $w_t(v)\!=\!\sum_{e\ni v}w_t(e)$ the total weight incident on $v$, the
regulator-aligned stressed set at horizon $h$ is
\begin{equation}
\Delta_t^h(v)\!=\!w_t(v)\!-\!w_{t+h}(v),\quad
\mathcal{S}_t^h \!=\! \mathrm{top}_{\pi}\bigl\{ v : w_t(v)\!>\!0,\;
                                              \Delta_t^h(v)\!>\!0 \bigr\},
\label{eq:stressed-app}
\end{equation}
with $\pi\!=\!0.05$ and ground-truth horizon $h\!=\!4$~weeks. The final
four test weeks lack the $4$-week-ahead snapshot $w_{t+4}$ and are not
scored, leaving $n\!=\!29$ evaluated weeks. This is the single ground-truth
definition used by \texttt{b2}, \texttt{b5}, and the LLM-decision pipeline.
A submission produces one JSON file per \{benchmark,~method\} pair plus a
single \texttt{results.json}; the harness fixes random seeds, snapshots
library versions, and emits SHA-256 hashes of all checkpoint files.

\subsection{Reference methods}
\label{app:reference-methods}
\begin{center}\small
\begin{tabularx}{\linewidth}{@{}lllX@{}}
\toprule
ID                              & Predictor      & Policy                & Role \\
\midrule
\texttt{h1\_weighted\_degree}   & ---            & heuristic alert       & Streaming anomaly baseline (\texttt{b2}) \\
\texttt{m1\_persistence\_rules} & $G_{t+h}\!=\!G_t$ & rule engine        & Decision baseline \\
\texttt{m2\_snapshot\_llm}      & current snap   & LLM                   & LLM-only ablation (no forecast) \\
\texttt{m3\_evolvegcn}          & EvolveGCN      & ---                   & Learned-GNN baseline \\
\texttt{m4\_fm\_only}           & \fmname{}      & ---                   & FM-only ablation \\
\texttt{m5\_fm\_rules}          & \fmname{}      & rule engine           & FM + rules \\
\texttt{m6\_fm\_llm}            & \fmname{}      & LLM                   & Full FM + LLM stack \\
\texttt{m7\_fm\_llm\_gated}     & \fmname{}      & LLM + safety gate     & Safety-gated FM + LLM \\
\bottomrule
\end{tabularx}
\end{center}
These methods expose all four cuts of the four-layer ablation: adding the
predictor (\texttt{m1}$\!\to\!$\texttt{m5}), swapping the decision layer
(\texttt{m5}$\!\to\!$\texttt{m6}), adding the safety gate
(\texttt{m6}$\!\to\!$\texttt{m7}), and removing the predictor from the LLM
stack (\texttt{m2}$\!\leftrightarrow\!$\texttt{m6}). All conform to a
shared \texttt{predict\_graph}\,/\,\texttt{decide\_actions} interface.

\paragraph{LLM model panel and decoding.}
All LLM-bearing methods (\texttt{m2}, \texttt{m6}, \texttt{m7}) are
evaluated under three decision-layer configurations: Claude~Opus~4.7
(main), Claude~Sonnet~4.6 (within-family, tier-below ablation), and
Google~Gemini~2.5~Pro (cross-family ablation). Each runs at temperature $0$
and $\max\_tokens=4096$; self-consistency is probed with three repeated
calls per week and Jaccard overlap on emitted ticket targets and
severities. The LLM-as-judge for the explanation-quality axis is a
three-model panel chosen so the judge tier is at or above the decision
tier: Claude~Opus~4.8 (strictly tier-above Opus~4.7),
Google~Gemini~2.5~Pro, and OpenAI~GPT-5.5; the Gemini and GPT judges run
with OpenRouter's $\mathrm{reasoning}=\{\mathrm{effort:minimal}\}$ so the
judge thinking budget matches the non-reasoning Anthropic baseline.
Claude~Haiku~4.5 is reported as a weak-judge reference. All models are
served through OpenRouter; observed cost is on the order of a few US
dollars per full sweep of the 29-week test split.
